Observed-range labels improve first-round balanced sampling by roughly six
percentage points, but a second acquired batch does not sustain that gain
against its matched continuation. The follow-up recovers a similar gain by
relabeling H0 alone. This identifies conditioning and its induced class
weighting as important parts of the result; adding configurations is not a
uniform benefit. Most of the gain is in HIGH-class fidelity, rather than
low-energy discovery.

The tested CVAEs do not establish a primary sampling advantage over the
no-latent autoregressive model. Their deterministic posterior means retain
a recoverable vacancy-ordering fingerprint. A fixed-condition random-encoder
control also contains this fingerprint, so recovery alone does not show
that training discovered it. That information does not become reliable
control of generated structures. In the illustrated fixed condition, tiny
mean variations coexist with nearly unit-scale posterior noise; a readable
mean scatter and a useful sampled latent are different outcomes.

Entropy estimation needs a different sampling objective. On the complete
12-site lattice, correct orbit multiplicities and proposal probabilities
are essential. Training toward a 1,000\,K Boltzmann target reduces the
equal-composition entropy MAE from 0.01310 to 0.00420\,$\kb$/cell for UMA
and from 0.01371 to 0.00313 for MACE. The same proposals are worse than
uniform sampling at 5,000\,K. Histogram flattening also does not establish
energy-bin convergence in the Wang--Landau calculation.

These findings concern fixed O3-derived vacancy lattices and the supplied
MLIP energies. Agreement across teachers is useful sensitivity evidence;
it does not validate DFT accuracy, relaxed phase stability, bulk entropy,
or molecular-dynamics predictions.
